\documentclass[11pt]{article}

\usepackage[margin=0.85in]{geometry}
\usepackage{amsmath,amssymb,bm,mathtools}
\usepackage{array,booktabs,longtable}
\usepackage{enumitem}
\usepackage{natbib}
\usepackage[hidelinks]{hyperref}
\usepackage{xcolor}

\newcommand{\R}{\mathbb{R}}
\newcommand{\Ts}{T_s}
\newcommand{\Tz}{T_z}
\newcommand{\mdot}{\dot m}
\newcommand{\QAC}{Q_{\mathrm{AC}}}
\newcommand{\QHX}{Q_{\mathrm{HVAC},X}}
\newcommand{\PHVAC}{P_{\mathrm{HVAC}}}
\newcommand{\dt}{\Delta t}
\newcommand{\PhiM}{\Phi_{\Delta t}}
\newcommand{\dd}{d}
\newcommand{\E}{\mathbb{E}}

\title{Part 6: Continuous-Time MPC, Observers, Forecasts, and Price-Aware Closed-Loop Control\\
Detailed Mathematical and Implementation Contract for PINODE/EPSR}
\author{}
\date{August 25, 2026}

\begin{document}
\maketitle

\section*{Purpose and relationship to Parts 1--5}

This document is the \textbf{Day-7 mathematical contract} for embedding the four already-defined
PINODE/EPSR thermal-model methods inside a continuous nonlinear model predictive controller (MPC).
It does not redefine Parts 1--5.  Instead, it specifies how the already-frozen continuous-time
thermal models, Phase-C HVAC surrogates, observers, forecasts, actuator limits, and EnergyPlus
closed-loop plant are connected.

The control formulation follows the standard receding-horizon MPC principle
\citep{garcia1989mpc,rawlings2026mpc,drgona2020all,killian2016ten}, while preserving the
continuous-time building-model structure used in grey-box building thermal modeling
\citep{bacher2011heat} and the paper's Neural-ODE/PINN lineage
\citep{chen2018node,raissi2019pinn}.

\subsection*{Locked non-negotiable contracts}

\begin{align}
\boxed{\text{All four paper methods remain continuous-time models.}}
\end{align}

\begin{align}
\boxed{\text{Every MPC state transition is produced by explicit numerical integration.}}
\end{align}

\begin{align}
\boxed{\text{No binary heating/cooling decision variable is introduced into the MPC.}}
\end{align}

\begin{align}
\boxed{\text{The final MPC uses }(\mdot,\Ts)\text{ as continuous physical control variables.}}
\end{align}

\begin{align}
\boxed{\QHX\rightarrow G_Q\rightarrow\QAC\rightarrow f_\theta
\quad\text{and}\quad
|\QAC|\rightarrow G_P\rightarrow\PHVAC.}
\end{align}

\begin{align}
\boxed{\PHVAC\text{ is an electrical objective/output quantity and never enters the thermal ODE RHS.}}
\end{align}

\begin{align}
\boxed{\text{TRAIN-supported actuator boxes are used in the paper experiments.}}
\end{align}

\begin{align}
\boxed{\text{Perfect and parameterized-noisy disturbance forecasts are both first-class modes.}}
\end{align}

\begin{align}
\boxed{\text{2C latent states are estimated online by observers.}}
\end{align}

\begin{align}
\boxed{\text{All Day-7 implementation remains inside }\texttt{Paper\_PINODE\_EPSR/src/pinode\_epsr}.}
\end{align}

\section{Sampled control time and continuous-time prediction}

The canonical sampled interval inherited from Parts 1--5 is
\begin{align}
\boxed{\dt=300\ {\rm s}.}
\end{align}
At control step \(k\), the physical time is
\begin{align}
t_k=t_0+k\dt.
\end{align}

The thermal model is never replaced by a learned discrete map.  Instead,
\begin{align}
\boxed{\dot x(t)=f_\theta\!\left(x(t),\QAC(t),d(t)\right)}
\tag{6.1}
\end{align}
is integrated over one sampled interval to form
\begin{align}
\boxed{
x_{k+1}
=
\PhiM\!\left(x_k,{\QAC}_k,d_k\right).
}
\tag{6.2}
\end{align}
Here \(\PhiM\) is a numerical integration operator and not an independently identified
discrete-time model.

The use of a continuous-time model inside MPC is especially natural for RC models and Neural
ODE models, where the governing object is the vector field itself
\citep{bacher2011heat,chen2018node}.

\section{Thermal forcing ownership}

For each thermal zone \(i\), define the non-HVAC thermal disturbances
\begin{align}
d_i(t)
=
\left[
T_o,\,
Q_{ZIC,i},\,
Q_{ZIR,i},\,
Q_{Sol1,i},\,
Q_{Sol2,i},\,
\text{architecture-specific coupling/disturbance channels}
\right].
\end{align}

The locked heat-channel decomposition is
\begin{align}
Q_{c,i}
&=
{\QAC}_i
+
Q_{ZIC,i}
+
Q_{Sol1,i},
\tag{6.3}\\
Q_{r,i}
&=
Q_{ZIR,i}
+
Q_{Sol2,i}.
\tag{6.4}
\end{align}

The terms \(Q_{ZIC}\), \(Q_{ZIR}\), \(Q_{Sol1}\), and \(Q_{Sol2}\) already contain the
appropriate thermal consequences of occupancy, lighting, plug/electric equipment, and solar
loading as defined upstream.  They remain \textbf{exogenous thermal disturbances}, not MPC
decision variables.  Building MPC literature likewise treats occupancy, weather, and solar/internal
loads as disturbances and forecasts rather than as HVAC actuation variables
\citep{oldewurtel2013occupancy,killian2016ten,drgona2020all}.

\section{RC dynamics used by MPC}

\subsection{1R1C}

For zone \(i\), the 1R1C state is
\begin{align}
x_i=T_i.
\end{align}
The continuous-time dynamics are
\begin{align}
\boxed{
C_i\dot T_i
=
\frac{T_o-T_i}{R_{io}}
+
Q_{c,i}
+
\eta_{r,i}Q_{r,i}
+
Q_{\mathrm{coupling},i}.
}
\tag{6.5}
\end{align}

For the paper's identifiable 1C construction, the radiative allocation is inherited unchanged
from Part 1.  MPC does not introduce a new RC parameterization.

\subsection{2R2C}

For zone \(i\),
\begin{align}
x_i=
\begin{bmatrix}
T_{a,i}\\
T_{m,i}
\end{bmatrix},
\end{align}
and
\begin{align}
\boxed{
C_{a,i}\dot T_{a,i}
=
\frac{T_o-T_{a,i}}{R_{io}}
+
\frac{T_{m,i}-T_{a,i}}{R_{im}}
+
Q_{c,i}
+
(1-\eta_i)Q_{r,i}
+
Q_{\mathrm{coupling},i},
}
\tag{6.6}
\end{align}
\begin{align}
\boxed{
C_{m,i}\dot T_{m,i}
=
\frac{T_{a,i}-T_{m,i}}{R_{im}}
+
\eta_iQ_{r,i}.
}
\tag{6.7}
\end{align}

All-to-One, Identity/IND, Identity/DEP1, and Identity/DEP2 retain the exact state ownership,
inter-zone coupling, and disturbance-allocation definitions established in Part 1.  In
particular, DEP1 retains reciprocal shared thermal coupling and DEP2 retains the additional
aggregate disturbance-allocation structure.  Part 6 only provides the MPC wrapper around those
already-defined vector fields.

The choice of reduced-order RC models as control-oriented building models is well established
in the building-MPC literature \citep{bacher2011heat,ma2012predictive,afram2014review}.

\section{Method-specific continuous-time MPC predictors}

Let \(a\) denote the Phase-D architecture and \(r\in\{1C,2C\}\) the state order.

\subsection{Inverse PINN-RC}

The trajectory PINN is used only for identification.  The deployed MPC model is the identified
physical RC ODE:
\begin{align}
\boxed{
\dot x
=
f_{\mathrm{RC}}
\left(
x,\QAC,d;\widehat\theta_{\mathrm{RC}}
\right).
}
\tag{6.8}
\end{align}
The identified parameters \(\widehat\theta_{\mathrm{RC}}\) are frozen before closed-loop MPC.
The inverse-PINN principle follows the broader physics-informed inverse-problem lineage
\citep{raissi2019pinn}.

\subsection{Neural ODE}

The normalized state is
\begin{align}
x=\mu_x+S_xz,
\end{align}
and normalized time is
\begin{align}
\boxed{
\tau=\frac{t-t_{\mathrm{ref}}}{\dt}.
}
\tag{6.9}
\end{align}
The learned vector field is
\begin{align}
\boxed{
\frac{dz}{d\tau}
=
g_\omega(z,\widetilde v),
}
\tag{6.10}
\end{align}
which corresponds to the physical-time derivative
\begin{align}
\boxed{
\dot x
=
\frac{S_x}{\dt}
g_\omega(z,\widetilde v).
}
\tag{6.11}
\end{align}
The NODE contains no RC residual and no energy projection
\citep{chen2018node}.

\subsection{Base PINODE}

The Base PINODE uses the same normalized-time neural derivative,
\begin{align}
\boxed{
\frac{dz}{d\tau}=g_\omega(z,\widetilde v),
}
\tag{6.12}
\end{align}
and therefore
\begin{align}
\boxed{
\widetilde f_\omega
=
\frac{S_x}{\dt}g_\omega.
}
\tag{6.13}
\end{align}
The RC equations influenced training through a soft residual.  At MPC deployment, the frozen
learned vector field is propagated exactly as defined in Part 4; no new hard projection is added.

\subsection{EBP-PINODE}

At each vector-field evaluation, the network first provides the raw physical derivative
\begin{align}
\widetilde f_\omega.
\end{align}
The exact energy-balance projection is
\begin{align}
f_P
=
\arg\min_f
\frac12
(f-\widetilde f_\omega)^\top
W
(f-\widetilde f_\omega)
\quad
\text{s.t.}
\quad
Af=b.
\tag{6.14}
\end{align}
Define
\begin{align}
\rho&=A\widetilde f_\omega-b,
\\
M&=AW^{-1}A^\top.
\end{align}
The multiplier is obtained by solving
\begin{align}
\boxed{
M\nu=\rho,
}
\tag{6.15}
\end{align}
and the projected derivative is
\begin{align}
\boxed{
f_P
=
\widetilde f_\omega
-
W^{-1}A^\top\nu.
}
\tag{6.16}
\end{align}
No explicit matrix inverse of \(M\) is introduced.  For the normalized solver,
\begin{align}
\boxed{
g_P
=
\dt S_x^{-1}f_P.
}
\tag{6.17}
\end{align}
The projection is evaluated at \textbf{every RK4 stage}, because \(A\) and \(b\) may depend on
the current stage state.

\section{Fixed-step RK4 contract}

Let \(N_s\ge1\) be the number of numerical substeps per 5-minute control interval.

\subsection{Physical-time form}

For a physical-time vector field
\begin{align}
\dot x=f(x,u,d),
\end{align}
define
\begin{align}
h_t=\frac{\dt}{N_s}.
\end{align}
For substep \(n\),
\begin{align}
k_1&=f(x_n,u_k,d_k),
\\
k_2&=f\!\left(x_n+\frac{h_t}{2}k_1,u_k,d_k\right),
\\
k_3&=f\!\left(x_n+\frac{h_t}{2}k_2,u_k,d_k\right),
\\
k_4&=f\!\left(x_n+h_tk_3,u_k,d_k\right),
\\
x_{n+1}
&=
x_n
+
\frac{h_t}{6}
(k_1+2k_2+2k_3+k_4).
\tag{6.18}
\end{align}

\subsection{Normalized-time form}

For NODE/Base-PINODE/EBP-PINODE, one sampled interval has
\begin{align}
\Delta\tau=1,
\end{align}
so the normalized substep is
\begin{align}
h_\tau=\frac{1}{N_s}.
\end{align}
The same RK4 algebra is applied to \(g_\omega\) or \(g_P\).

\subsection{Zero-order-hold forcing}

Inside one 5-minute prediction interval,
\begin{align}
u(\tau)&=u_k,
\\
d(\tau)&=d_k,
\qquad
\tau\in[\tau_k,\tau_k+1),
\tag{6.19}
\end{align}
unless a later explicitly documented experiment introduces an intra-interval forcing model.
This preserves the Parts 3--5 numerical convention.

\subsection{CasADi parity requirement}

The CasADi expression for every frozen model must satisfy a direct numerical parity test against
the authoritative PyTorch/Neuromancer implementation:
\begin{align}
\boxed{
\left\|
\Phi_{\mathrm{Torch}}(x,u,d)
-
\Phi_{\mathrm{CasADi}}(x,u,d)
\right\|
\le\epsilon_{\mathrm{parity}}.
}
\tag{6.20}
\end{align}
CasADi is used as a symbolic/algorithmic differentiation and nonlinear optimal-control framework
\citep{andersson2019casadi}; it does not redefine the learned physics.

\section{HVAC algebraic chain inside MPC}

\subsection{Physical HVAC proxy}

Let
\begin{align}
c_{p,a}=1000(1.005)=1005\ {\rm J/(kg\,K)}.
\end{align}
Then the Phase-C HVAC proxy is
\begin{align}
\boxed{
\QHX
=
c_{p,a}\mdot(\Ts-\Tz).
}
\tag{6.21}
\end{align}

\subsection{Why \(G_Q\) remains necessary}

The thermal-control signal used by the learned thermal models is
\begin{align}
\boxed{
\QAC
=
G_Q(\QHX).
}
\tag{6.22}
\end{align}
For the current controlled paper case, \(G_Q\) is a persisted Phase-C one-input surrogate,
typically affine:
\begin{align}
G_Q(q)=a_Qq+b_Q.
\tag{6.23}
\end{align}

Although \(\QHX\) can be computed directly from \((\mdot,\Ts,\Tz)\), it is not generally
identical to the Phase-D \(\QAC\) signal on which the thermal models were trained.  Therefore
bypassing \(G_Q\) would change the learned model's input semantics after training.

\subsection{HVAC electric power}

The common Phase-C electric-power surrogate is
\begin{align}
\boxed{
P_{\mathrm{raw}}
=
G_P(q^{\mathrm{abs}}),
}
\tag{6.24}
\end{align}
with \(G_P\) typically affine:
\begin{align}
G_P(q)=a_Pq+b_P.
\tag{6.25}
\end{align}
The realized electrical quantity used for energy accounting is constrained to be nonnegative.

\section{Smooth continuous epigraph for \texorpdfstring{\(|\QAC|\)}{|QAC|} and nonnegative power}

To avoid introducing a nonsmooth absolute-value operator into IPOPT, introduce the continuous
auxiliary variable
\begin{align}
q^{\mathrm{abs}}_j\ge0
\end{align}
and impose
\begin{align}
q^{\mathrm{abs}}_j&\ge {\QAC}_j,
\\
q^{\mathrm{abs}}_j&\ge-{\QAC}_j.
\tag{6.26}
\end{align}
Introduce
\begin{align}
P_j\ge0
\end{align}
and impose
\begin{align}
P_j\ge G_P(q^{\mathrm{abs}}_j).
\tag{6.27}
\end{align}
Provided the validated \(G_P\) has nonnegative slope and \(P_j\) carries a positive objective
weight, minimization makes the epigraph tight:
\begin{align}
q^{\mathrm{abs}}_j&=|{\QAC}_j|,
\\
P_j&=\max\!\left(0,G_P(|{\QAC}_j|)\right).
\tag{6.28}
\end{align}
Thus the NLP remains continuous and compatible with the smooth nonlinear-programming framework
used by IPOPT \citep{wachter2006ipopt}.

\section{Frozen continuous HVAC control architecture}

Day-7 actuator forensics and the direct EnergyPlus Runtime-API implementation
close the earlier actuator-architecture branch.  The final paper architecture
uses both mass flow and supply-air temperature as continuous physical MPC
variables:
\begin{align}
\boxed{
u_{i,j}
=
\begin{bmatrix}
\mdot_{i,j}\\
{\Ts}_{i,j}
\end{bmatrix},
\qquad
i\in\{D,K\}.
}
\tag{6.29}
\end{align}

For the identity two-zone plant the implemented first control move is therefore
the four-command vector
\begin{align}
\boxed{
u_k^\star
=
\begin{bmatrix}
\mdot_{D,k}^\star &
{\Ts}_{D,k}^\star &
\mdot_{K,k}^\star &
{\Ts}_{K,k}^\star
\end{bmatrix}^{\!\top}.
}
\tag{6.30}
\end{align}

Heating and cooling emerge continuously from the sign of the air-side thermal
delivery,
\begin{align}
{\QHX}_{i,j}
=
\mdot_{i,j}c_p({\Ts}_{i,j}-{\Tz}_{i,j}),
\end{align}
so no binary heating/cooling decision variable is introduced.

The low-level EnergyPlus realization is identical for Dining and Kitchen:
\begin{align}
a_{m,i,j}
&=
\tfrac12\mdot_{i,j},
\\
a_{q,i,j}
&=
\mdot_{i,j}c_p({\Ts}_{i,j}-{\Tz}_{i,j}).
\tag{6.31}
\end{align}
The fan command is issued inside the HVAC system iteration loop and the unitary
sensible-load request is issued after the predictor/HVAC managers, as frozen by
Part~7.

The earlier \(\mdot\)-only alternative is retired from the final paper
architecture.  It may remain in historical development notes, but it is not a
second experimental actuator path and is not a factor in the final
model-to-model comparison.

\subsection{Generic per-zone feasibility supervisor}

The simulator applies one symmetric numerical feasibility test to Dining and
Kitchen and to both signs of \({\Ts}-{\Tz}\).  The current configured
Day-7 envelope is
\begin{align}
0.50
\le
\frac{\mdot_i^\star}{\mdot_{i,\max}}
\le
0.80,
\qquad
4^\circ{\rm C}
\le
|{\Ts}_i^\star-{\Tz}_i|
\le
8^\circ{\rm C}.
\tag{6.32}
\end{align}

All four regimes---Dining heating, Dining cooling, Kitchen heating, and Kitchen
cooling---are MPC-commandable.  There is no categorical Kitchen-heating
exclusion.  If the command for zone \(i\) is outside the configured envelope,
both external overrides for that zone are released for the current 300-s
interval and native EnergyPlus HVAC control is used for that zone only.  The
other zone remains under MPC control if its command is feasible.

The exact final optimization bounds may be tightened after learned-model
training so that the MPC also respects model TRAIN support.  This does not
change the simulator supervisor mechanism.

\section{TRAIN-supported actuator boxes}

For each zone \(i\) and observed HVAC regime \(m\), define the exact TRAIN+included active-mode
sample set
\begin{align}
\mathcal A_{i,m}^{\mathrm{tr}}
=
\left\{
(\mdot_k,{\Ts}_k):
k\in\mathrm{TRAIN+included},\;
\mathrm{mode}_k=m
\right\}.
\tag{6.35}
\end{align}

Store at least
\begin{align}
\min,\quad
\max,\quad
\mathrm{mean},\quad
\mathrm{median},\quad
\mathrm{std},\quad
P_{05},P_{10},P_{25},P_{50},P_{75},P_{90},P_{95}
\tag{6.36}
\end{align}
for both \(\mdot\) and \(\Ts\).

The primary paper control boxes are data-supported:
\begin{align}
\boxed{
\mdot_{i,m}^{\min,\mathrm{tr}}
\le
\mdot_{i,j}
\le
\mdot_{i,m}^{\max,\mathrm{tr}},
}
\tag{6.37}
\end{align}
and
\begin{align}
\boxed{
{\Ts}_{i,m}^{\min,\mathrm{tr}}
\le
{\Ts}_{i,j}
\le
{\Ts}_{i,m}^{\max,\mathrm{tr}}.
}
\tag{6.38}
\end{align}

EnergyPlus hard actuator limits are recorded separately as plant constraints.  The paper
experiments use the intersection
\begin{align}
\boxed{
\mathcal U_{\mathrm{paper}}
=
\mathcal U_{\mathrm{TRAIN}}
\cap
\mathcal U_{\mathrm{act}},
}
\tag{6.39}
\end{align}

The simulator permits all four zone/mode combinations, but the predictive
model must still be used responsibly.  If a final learned-model checkpoint has
weak or absent TRAIN support for a requested region, that command is marked
model-OOD and should not be part of the primary fair-comparison operating
region unless explicitly studied as robustness/extrapolation.  This
model-support rule is distinct from the generic EnergyPlus native-fallback
supervisor.

\section{Condensed All-to-One physical-input allocation}

The All-to-One learned model predicts one aggregate thermal state, while the
EnergyPlus plant still contains two physical PSZ systems.  Therefore the final
architecture contains an explicit aggregate-to-zone allocation layer before
the common actuator transform.

Let
\begin{align}
u_{A,k}^{\star}
=
\begin{bmatrix}
\mdot_{A,k}^{\star}\\
{\Ts}_{A,k}^{\star}
\end{bmatrix}.
\tag{6.40}
\end{align}
The physical zone commands are
\begin{align}
\boxed{
u_{k}^{z,\star}
=
\mathcal A_u
\left(
u_{A,k}^{\star},
\xi_{D,k},
\xi_{K,k}
\right)
=
\begin{bmatrix}
u_{D,k}^{\star}\\
u_{K,k}^{\star}
\end{bmatrix},
}
\tag{6.41}
\end{align}
where each \(u_{i,k}^{\star}=[\mdot_{i,k}^\star,{\Ts}_{i,k}^\star]^\top\).

The allocator is required to preserve the aggregate physical-input meaning
defined upstream.  With aggregation map \(P_u\), record
\begin{align}
\boxed{
r_{u,k}
=
u_{A,k}^{\star}
-
P_u u_{k}^{z,\star}
}
\tag{6.42}
\end{align}
as an allocation-consistency diagnostic.

The allocation rule is common across learned-model families and is frozen
before final comparative testing.  After allocation, Dining and Kitchen pass
through the same feasibility supervisor and the same low-level actuator
transformation.  The aggregate predictive model therefore does not bypass or
merge the two real EnergyPlus HVAC systems.

\section{Disturbance forecasts}

Building MPC obtains much of its value from using forecasts of weather and internal disturbances
\citep{killian2016ten,oldewurtel2014stochastic,drgona2020all}.

Let
\begin{align}
d_{k+j}
\end{align}
be the true exogenous non-HVAC disturbance vector over the horizon.

\subsection{Perfect forecast}

The perfect-information benchmark is
\begin{align}
\boxed{
\widehat d_{k+j|k}
=
d_{k+j},
\qquad
j=0,\ldots,N_p-1.
}
\tag{6.42}
\end{align}
Future measured \(\Tz\) and future historical \(\QAC\) are never supplied as forecasts.

\subsection{Parameterized noisy forecast}

For disturbance component \(r\),
\begin{align}
\widehat d^{(r)}_{k+j|k}
=
d^{(r)}_{k+j}
+
b_r
+
\epsilon^{(r)}_{k+j|k},
\tag{6.43}
\end{align}
where
\begin{align}
\epsilon^{(r)}_{k+j|k}
=
\rho_r\epsilon^{(r)}_{k+j-1|k}
+
\sqrt{1-\rho_r^2}\,
\sigma^{(r)}_j
z^{(r)}_{k+j},
\qquad
z^{(r)}_{k+j}\sim\mathcal N(0,1).
\tag{6.44}
\end{align}
Use
\begin{align}
\sigma^{(r)}_j
=
\gamma_r(j)
\left(
\sigma^{(r)}_{\mathrm{abs}}
+
\sigma^{(r)}_{\mathrm{rel}}
|d^{(r)}_{k+j}|
\right),
\tag{6.45}
\end{align}
where \(\gamma_r(j)\ge1\) optionally increases uncertainty with forecast lead time.

The forecast configuration therefore exposes
\begin{align}
\boxed{
b_r,\;
\sigma^{(r)}_{\mathrm{abs}},\;
\sigma^{(r)}_{\mathrm{rel}},\;
\rho_r,\;
\gamma_r(\cdot),\;
\mathrm{seed}.
}
\tag{6.46}
\end{align}
Setting all error terms to zero exactly recovers the perfect forecast.

\section{Observer contract}

\subsection{1C models}

For a 1C model, the complete thermal state is the measured EnergyPlus zone-air temperature:
\begin{align}
\boxed{
\widehat x_k
=
T_{z,k}^{\mathrm{EP}}.
}
\tag{6.47}
\end{align}
No observer is required.

\subsection{Learned 2C methods: causal encoder as online observer}

For NODE-2C, Base-PINODE-2C, and EBP-PINODE-2C, let the rolling causal history be
\begin{align}
\mathcal H_k
=
\left\{
T_{a,\ell}^{\mathrm{meas}},
{\QAC}_{\ell}^{\mathrm{realized}},
d_\ell
\right\}_{\ell=k-L_e+1}^{k}.
\tag{6.48}
\end{align}
The frozen trained causal encoder is re-evaluated at every \textbf{closed-loop control boundary}:
\begin{align}
\boxed{
\widehat T_{m,k}
=
E_{\psi^\star}(\mathcal H_k),
}
\tag{6.49}
\end{align}
and
\begin{align}
\boxed{
\widehat x_k
=
\begin{bmatrix}
T_{a,k}^{\mathrm{EP}}\\
\widehat T_{m,k}
\end{bmatrix}.
}
\tag{6.50}
\end{align}

Inside one MPC prediction horizon, the encoder is \textbf{not} repeatedly called on predicted
future measurements.  The horizon starts from \(\widehat x_k\) and then propagates only through
the continuous-time model and numerical integrator.

The realized historical control input supplied to the encoder is constructed causally from
realized/commanded plant variables:
\begin{align}
{\QAC}_{\ell}^{\mathrm{realized}}
=
G_Q
\left[
c_{p,a}
\mdot_{\ell}^{\mathrm{realized}}
\left(
{\Ts}_{\ell}^{\mathrm{realized}}
-
T_{z,\ell}^{\mathrm{meas}}
\right)
\right],
\tag{6.51}
\end{align}
unless a directly measured authoritative plant signal is explicitly substituted with provenance.

\subsection{Inverse PINN-RC 2C: Luenberger observer}

For the identified RC model with output
\begin{align}
y=Hx,
\qquad
H=
\begin{bmatrix}
1&0
\end{bmatrix},
\tag{6.52}
\end{align}
a continuous Luenberger-style observer is
\begin{align}
\boxed{
\dot{\widehat x}
=
f_{\mathrm{RC}}(\widehat x,\QAC,d)
+
L
\left(
y-H\widehat x
\right).
}
\tag{6.53}
\end{align}
For a linear RC realization
\begin{align}
\dot x=Ax+B_u u+B_dd,
\end{align}
the estimation error satisfies
\begin{align}
\dot e=(A-LH)e,
\tag{6.54}
\end{align}
so \(L\) is chosen to place the observer error poles in the stable half-plane, following the
classical observer construction of \citet{luenberger1966observer}.

A sampled implementation may equivalently use
\begin{align}
\widehat x_{k|k-1}
&=
\Phi_{\mathrm{RC},\Delta t}
(
\widehat x_{k-1|k-1},{\QAC}_{k-1},d_{k-1}
),
\\
\widehat x_{k|k}
&=
\widehat x_{k|k-1}
+
L_d
\left(
y_k-H\widehat x_{k|k-1}
\right).
\tag{6.55}
\end{align}

\subsection{Inverse PINN-RC 2C: Kalman/Kalman-family observer}

The second supported observer is covariance based.  The state prediction uses the \textbf{same
numerically integrated RC transition}:
\begin{align}
\widehat x^-_k
=
\Phi_{\mathrm{RC},\Delta t}
(
\widehat x^+_{k-1},{\QAC}_{k-1},d_{k-1}
).
\tag{6.56}
\end{align}
Let
\begin{align}
F_{k-1}
=
\left.
\frac{\partial
\Phi_{\mathrm{RC},\Delta t}(x,{\QAC}_{k-1},d_{k-1})}
{\partial x}
\right|_{x=\widehat x^+_{k-1}}.
\tag{6.57}
\end{align}
Then
\begin{align}
P^-_k
&=
F_{k-1}P^+_{k-1}F_{k-1}^\top+Q_o,
\\
K_k
&=
P^-_kH^\top
\left(
HP^-_kH^\top+R_o
\right)^{-1},
\\
\widehat x^+_k
&=
\widehat x^-_k
+
K_k
\left(
y_k-H\widehat x^-_k
\right),
\\
P^+_k
&=
(I-K_kH)P^-_k.
\tag{6.58}
\end{align}
For the linear RC case this is the standard Kalman filter; the Jacobian form also supports an
extended-Kalman implementation if required.  The lineage is the classical state-estimation
framework of \citet{kalman1960filter}.

\section{Electricity price signal}

Let \(\pi_k\) denote the electricity-price signal used by the economic MPC.

The primary EPSR experiment uses the already-established daily normalized price profile from the
MS-thesis workflow.  That profile is treated as experiment input data, not as a newly invented
utility tariff.  If the stored profile has 10-minute resolution and the MPC interval is 5 minutes,
use zero-order hold:
\begin{align}
\boxed{
\pi_{2r}=\pi_{2r+1}=\pi_r^{10\mathrm{min}}.
}
\tag{6.59}
\end{align}

A normalized price produces a \textbf{normalized price-weighted energy metric}, not a monetary
dollar cost.  If a future experiment substitutes a real tariff in \$/kWh, the same equations
become a true monetary objective.

Price-responsive HVAC MPC is directly motivated by demand-response and dynamic-pricing
literature \citep{avci2013price,amin2020price,zanetti2025flex}.

\section{Why non-HVAC electric loads are not optimized}

Let
\begin{align}
P_{\mathrm{base},k}
=
P_{\mathrm{lights},k}
+
P_{\mathrm{equipment},k}
+
P_{\mathrm{other,nonHVAC},k}.
\tag{6.60}
\end{align}
These loads may be reconstructed from Phase-B schedule/rating information for reporting, but they
are \textbf{not controlled} in the present HVAC MPC.

If the objective contained only linear energy/price terms,
\begin{align}
\sum_j
\pi_j
\left(
{\PHVAC}_j
+
P_{\mathrm{base},j}
\right)\dt,
\end{align}
and \(P_{\mathrm{base},j}\) is exogenous, then
\begin{align}
\sum_j\pi_jP_{\mathrm{base},j}\dt
\end{align}
is constant with respect to the MPC decision variables and therefore does not change the optimizer.

Thus the main MPC should optimize \(\PHVAC\) only.  For EPSR reporting, however, define
\begin{align}
\boxed{
P_{\mathrm{building},k}
=
{\PHVAC}_k
+
P_{\mathrm{base},k},
}
\tag{6.61}
\end{align}
so total-building demand, grid-facing peaks, and total price-weighted consumption can be plotted
without claiming that lighting/equipment were controlled.

An exception would arise if a later formulation adds a whole-building peak-demand charge,
explicit controllable lighting/plug loads, or another coupling constraint.  In that case the
non-HVAC electrical profile can affect optimal HVAC timing and must be included explicitly.
That extension is outside the present Part-6 lock.

\section{Comfort model}

Let the desired setpoint and deadband/comfort width at horizon point \(j\) be
\begin{align}
T^{\mathrm{sp}}_j,
\qquad
\delta^{\mathrm{db}}_j\ge0.
\end{align}
Define
\begin{align}
T^{\mathrm{low}}_j
&=
T^{\mathrm{sp}}_j-\delta^{\mathrm{db}}_j,
\\
T^{\mathrm{high}}_j
&=
T^{\mathrm{sp}}_j+\delta^{\mathrm{db}}_j.
\tag{6.62}
\end{align}
Introduce nonnegative comfort slacks
\begin{align}
s_j^-\ge0,
\qquad
s_j^+\ge0,
\end{align}
with
\begin{align}
T^{\mathrm{low}}_j-s_j^-
\le
T_{z,j}
\le
T^{\mathrm{high}}_j+s_j^+.
\tag{6.63}
\end{align}
Soft comfort constraints prevent infeasibility while making violations explicit
\citep{ma2012predictive,drgona2020all}.

\section{Price-aware MPC objective}

For the frozen two-input architecture, define
\begin{align}
J
=
\sum_{j=0}^{N_p-1}
\Big[
&
w_c
\left(
(s_j^-)^2+(s_j^+)^2
\right)
+
w_e
P_j\dt_h
+
w_p
\pi_jP_j\dt_h
\nonumber\\
&
+
w_m
(\mdot_j-\mdot_{j-1})^2
+
w_T
({\Ts}_j-{\Ts}_{j-1})^2
\Big]
+
J_f,
\tag{6.64}
\end{align}
where
\begin{align}
\dt_h=\frac{\dt}{3600}
\end{align}
and \(J_f\) is an optional terminal comfort term.

The terms have distinct purposes:
\begin{align}
J_{\mathrm{comfort}}
&\rightarrow
\text{maintain thermal comfort},
\\
J_{\mathrm{energy}}
&\rightarrow
\text{reduce HVAC electrical energy},
\\
J_{\mathrm{price}}
&\rightarrow
\text{shift HVAC electricity away from expensive periods},
\\
J_{\Delta u}
&\rightarrow
\text{avoid unnecessarily aggressive actuator movement}.
\end{align}

Keeping both \(w_eP_j\) and \(w_p\pi_jP_j\) is deliberate when \(\pi\) is a normalized synthetic
price signal: energy is still penalized even during low-price periods.  Dynamic-price HVAC MPC
has been widely studied as a building demand-response mechanism
\citep{avci2013price,amin2020price,zanetti2025flex}.

For the retired mass-flow-only alternative, remove the \(T_s\)-ramping term:
\begin{align}
w_T=0,
\qquad
{\Ts}_j\text{ known externally.}
\tag{6.65}
\end{align}

\section{Complete nonlinear program: the frozen two-input architecture}

At control time \(k\), the decision variables include
\begin{align}
\mathcal Z_k
=
\left\{
x_j,\mdot_j,{\Ts}_j,
q_j^{\mathrm{abs}},P_j,s_j^-,s_j^+
\right\}_{j=0}^{N_p}.
\tag{6.66}
\end{align}
Solve
\begin{align}
\boxed{
\min_{\mathcal Z_k}J
}
\tag{6.67}
\end{align}
subject to:

\paragraph{Observer-initialized state}
\begin{align}
x_0=\widehat x_k.
\tag{6.68}
\end{align}

\paragraph{HVAC proxy and Phase-C thermal input}
\begin{align}
{\QHX}_j
&=
c_{p,a}\mdot_j({\Ts}_j-T_{z,j}),
\\
{\QAC}_j
&=
G_Q({\QHX}_j).
\tag{6.69}
\end{align}

\paragraph{Continuous-time thermal prediction with numerical integration}
\begin{align}
x_{j+1}
=
\Phi_{\Delta t}
\left(
x_j,{\QAC}_j,\widehat d_{k+j|k};N_s
\right).
\tag{6.70}
\end{align}

\paragraph{Absolute-QAC epigraph}
\begin{align}
q_j^{\mathrm{abs}}&\ge{\QAC}_j,
\\
q_j^{\mathrm{abs}}&\ge-{\QAC}_j,
\\
q_j^{\mathrm{abs}}&\ge0.
\tag{6.71}
\end{align}

\paragraph{PHVAC epigraph}
\begin{align}
P_j&\ge G_P(q_j^{\mathrm{abs}}),
\\
P_j&\ge0.
\tag{6.72}
\end{align}

\paragraph{Comfort}
\begin{align}
T^{\mathrm{low}}_{k+j}-s_j^-
\le
T_{z,j}
\le
T^{\mathrm{high}}_{k+j}+s_j^+,
\\
s_j^\pm\ge0.
\tag{6.73}
\end{align}

\paragraph{Data-supported actuator box}
\begin{align}
\mdot^{\min,\mathrm{tr}}
&\le\mdot_j\le\mdot^{\max,\mathrm{tr}},
\\
{\Ts}^{\min,\mathrm{tr}}
&\le{\Ts}_j\le{\Ts}^{\max,\mathrm{tr}}.
\tag{6.74}
\end{align}

\paragraph{Optional single-regime sign constraint}
\begin{align}
{\QAC}_j\le0
\quad\text{or}\quad
{\QAC}_j\ge0.
\tag{6.75}
\end{align}

The resulting problem is a continuous nonlinear program suitable for CasADi and IPOPT
\citep{andersson2019casadi,wachter2006ipopt}.

\section{Complete nonlinear program: the retired mass-flow-only alternative}

the retired mass-flow-only alternative uses the same program except
\begin{align}
{\Ts}_j
=
{\Ts}_{m_j}^{\mathrm{fixed}}
\end{align}
is a known parameter and not a decision variable.  The control sequence becomes
\begin{align}
\mathcal U_k
=
\{\mdot_0,\ldots,\mdot_{N_p-1}\}.
\tag{6.76}
\end{align}
All equations for \(\QHX\), \(G_Q\), the continuous thermal model, RK4, \(G_P\), comfort, price,
and mass-flow ramping remain unchanged.

\section{Multi-zone identity cases}

For Dining/Kitchen identity architectures,
\begin{align}
u_j
=
\begin{bmatrix}
u_{D,j}\\
u_{K,j}
\end{bmatrix}.
\tag{6.77}
\end{align}
For the frozen two-input architecture,
\begin{align}
u_{i,j}
=
\begin{bmatrix}
\mdot_{i,j}\\
{\Ts}_{i,j}
\end{bmatrix};
\end{align}
for the retired mass-flow-only alternative,
\begin{align}
u_{i,j}=\mdot_{i,j}.
\end{align}

IND, DEP1, and DEP2 change the thermal model structure and coupling, not the physical ownership of
the Dining and Kitchen HVAC actuators.

\section{Receding-horizon closed-loop algorithm}

At each plant interval \(k\):

\begin{enumerate}[leftmargin=2em]
\item Read the current EnergyPlus measurements and realized actuator values.
\item Update the causal history \(\mathcal H_k\).
\item Estimate \(\widehat x_k\) with direct measurement (1C), causal encoder (learned 2C), or
Luenberger/Kalman observer (Inverse PINN-RC 2C).
\item Construct \(\widehat d_{k:k+N_p-1|k}\) using the selected forecast provider.
\item Load the known price horizon \(\pi_{k:k+N_p-1}\).
\item Solve the continuous NLP.
\item Apply only the first optimized control command \(u_k^\star\) to EnergyPlus.
\item Advance EnergyPlus by one control interval.
\item Record commanded and realized actuators, solver diagnostics, predicted trajectories, observed
temperatures, \(Q_{\mathrm{HVAC},X}\), \(\QAC\), \(\PHVAC\), comfort, and provenance.
\item Repeat.
\end{enumerate}

This is the standard receding-horizon principle \citep{garcia1989mpc,rawlings2026mpc}.

\section{Direct EnergyPlus Runtime-API plant boundary}

The final plant interface is a project-owned synchronous wrapper around the
EnergyPlus~24.1 Python Runtime/Data Exchange API.  the direct EnergyPlus wrapper and Gymnasium are
not dependencies of the paper closed-loop stack.

The interface is
\begin{verbatim}
observer / learned thermal model
              |
              v
             MPC
              |
              v
[m_dot_D*, Tsa_D*, m_dot_K*, Tsa_K*]
              |
              v
generic EnergyPlus simulator
  - feasibility supervisor
  - actuator transformation
  - callback execution
  - broad history recorder
              |
              v
       EnergyPlus 24.1
\end{verbatim}

The simulator records received physical commands, feasibility decisions,
transformed low-level commands, effective actuator readback, raw API actuator
values for forensics, zone/environment signals, the known PSZ air path,
component signals, derived air-side quantities, and both EnergyPlus
system-timestep and 300-s control-step histories.

Scientific run artifacts are written under the separate
\texttt{Data/ScaleBridge/Paper\_PINODE\_EPSR} root.  Clean source code, tests,
configuration, and mathematical contracts remain under the repository
\texttt{scalebridge-research/Paper\_PINODE\_EPSR} root.

\section{Price and disturbance knowledge in the MPC horizon}

For the primary perfect-information benchmark:
\begin{align}
\widehat d_{k+j|k}&=d_{k+j},
\\
\widehat \pi_{k+j|k}&=\pi_{k+j}.
\tag{6.78}
\end{align}
The perfect disturbance benchmark is useful because it isolates model/controller quality from
forecast error.  A second experiment introduces the parameterized noisy forecast of
Eqs.~(6.43)--(6.46), following the broader building-MPC uncertainty literature
\citep{oldewurtel2014stochastic,drgona2020all}.

\section{Recommended primary EPSR control experiment}

The primary MPC comparison should preferentially use a day/window satisfying:

\begin{align}
\boxed{\text{all controlled zones have direct TRAIN support for the active HVAC regime}}
\end{align}
and
\begin{align}
\boxed{\text{the selected actuator box remains inside the observed TRAIN envelope.}}
\end{align}

A cooling-dominated day remains an attractive primary comparison window
because it is strongly represented in the actuator-validation evidence and
typically provides simultaneous two-zone conditioning.  This is a
comparability choice, not a categorical restriction on Kitchen heating.

Heating windows may also be used.  All four zone/mode combinations are
MPC-commandable when the numerical supervisor accepts the command.  Separate
robustness tests should deliberately exercise per-zone native fallback.

\section{Evaluation metrics}

For each method, architecture, RC order, and forecast mode, save at least:

\subsection*{Comfort}
\begin{align}
\mathrm{RMSE}_{sp},
\quad
\mathrm{MAE}_{sp},
\quad
\mathrm{degree\ hours\ outside\ deadband},
\quad
\mathrm{fraction\ within\ deadband}.
\end{align}

\subsection*{HVAC energy}
\begin{align}
\boxed{
E_{\mathrm{HVAC}}
=
\sum_k
{\PHVAC}_k
\frac{\dt}{3600}
\frac{1}{1000}
\quad[\mathrm{kWh}].
}
\tag{6.79}
\end{align}

\subsection*{Normalized price-weighted HVAC energy}
For dimensionless thesis price \(\pi_k\),
\begin{align}
\boxed{
C_{\pi,\mathrm{HVAC}}^{\mathrm{norm}}
=
\sum_k
\pi_k
{\PHVAC}_k
\frac{\dt}{3600}
\frac{1}{1000}.
}
\tag{6.80}
\end{align}
This is \textbf{not} labeled dollars.

\subsection*{Grid-facing reporting}
\begin{align}
P_{\mathrm{building},k}
=
{\PHVAC}_k+P_{\mathrm{base},k},
\end{align}
with
\begin{align}
\max_k P_{\mathrm{building},k},
\qquad
\sum_k\pi_kP_{\mathrm{building},k}\dt_h/1000
\end{align}
reported optionally.

\subsection*{Actuation}
\begin{align}
\max|\Delta\mdot|,
\quad
\sum(\Delta\mdot)^2,
\quad
\max|\Delta\Ts|,
\quad
\sum(\Delta\Ts)^2,
\end{align}
as applicable.

\subsection*{Optimization}
\begin{align}
\mathrm{IPOPT\ success\ fraction},
\quad
\mathrm{solve\ time\ mean/median/P95/max},
\quad
\mathrm{iteration\ count},
\quad
\mathrm{constraint\ residuals}.
\end{align}

\subsection*{Forecast robustness}
\begin{align}
\mathrm{forecast\ RMSE},
\quad
\mathrm{comfort\ degradation},
\quad
\mathrm{energy/cost\ degradation}
\end{align}
versus perfect forecast.

\subsection*{Data-support/OOD}
Record fractions of commands outside empirical model-support quantiles, on TRAIN-envelope
boundaries, rejected by the configured actuator-feasibility supervisor, or marked model-OOD.

\section{Fair-comparison contract}

Across Inverse PINN-RC, NODE, Base PINODE, and EBP-PINODE:

\begin{align}
\boxed{\text{same plant episode}}
\end{align}
\begin{align}
\boxed{\text{same price trajectory}}
\end{align}
\begin{align}
\boxed{\text{same disturbance forecast realization and random seed}}
\end{align}
\begin{align}
\boxed{\text{same comfort bounds}}
\end{align}
\begin{align}
\boxed{\text{same four-command interface, supervisor, and data-supported bounds}}
\end{align}
\begin{align}
\boxed{\text{same MPC horizon and objective weights}}
\end{align}
unless a method-specific requirement is declared before seeing final test outcomes.

Model hyperparameters are frozen before final testing.  MPC weights and horizon are not retuned
against final comparative outcomes.

\section{Implementation invariants}

The code implementation must satisfy all of the following:

\begin{enumerate}[leftmargin=2em]
\item No learned method is converted into a discrete neural predictor.
\item Inverse PINN-RC is numerically integrated like the other continuous-time methods.
\item EBP projection is applied at every RK4 stage.
\item CasADi uses linear solves, not explicit projection-matrix inverses.
\item \(G_Q\) is retained between \(\QHX\) and the thermal-model \(\QAC\) input.
\item \(G_P\) receives the absolute-QAC epigraph quantity.
\item \(\PHVAC\) never enters the thermal model.
\item One frozen four-command HVAC architecture is used for all learned-model families.
\item No binary heat/cool variable appears in the NLP.
\item Learned 2C encoders are re-used as observers at real control boundaries.
\item The Inverse PINN-RC 2C controller supports both Luenberger and Kalman/Kalman-family observers.
\item Perfect and noisy forecasts share one forecast-provider interface.
\item Non-HVAC thermal gains remain disturbances.
\item Non-HVAC electrical loads are reporting quantities unless a later explicitly controllable
grid formulation is introduced.
\item All-to-One physical inputs pass through one recorded aggregate-to-zone allocator \(\mathcal A_u\).
\item The main paper actuator bounds are generated from exact TRAIN+included support and intersected
with the configured EnergyPlus actuator-feasibility envelope.
\end{enumerate}

\section{Scientific interpretation}

This Part-6 design places the paper in a useful intersection of:

\begin{align}
\boxed{\text{continuous-time control-oriented building modeling}}
\end{align}
\begin{align}
\boxed{\text{physics-informed / neural differential-equation models}}
\end{align}
\begin{align}
\boxed{\text{observer-based output feedback}}
\end{align}
\begin{align}
\boxed{\text{economic MPC under dynamic electricity prices}}
\end{align}
\begin{align}
\boxed{\text{building-grid demand flexibility}}
\end{align}

MPC for building thermal systems has a mature control literature
\citep{ma2012predictive,afram2014review,killian2016ten,drgona2020all}, and dynamic-price control
connects building thermal flexibility directly to grid-oriented demand response
\citep{avci2013price,amin2020price,zanetti2025flex}.  The paper's distinct contribution is not
the existence of MPC itself; it is the controlled comparison of the four continuous-time learned
thermal-model families, including the energy-balance-projected PINODE, inside the same
observer/forecast/actuator/economic-control contract.

\bibliographystyle{plainnat}
\bibliography{PINODE_EPSR_Part6_MPC_Observers_References}

\end{document}
